\documentclass[10pt,a4paper]{amsart}
\usepackage{amsmath,amssymb,amsthm,graphicx}
\usepackage{verbatim}
\newenvironment{reference}{\comment}{\endcomment}
%\newenvironment{reference}{}{}
\newenvironment{slogan}{\comment}{\endcomment}
\newenvironment{history}{\comment}{\endcomment}
% For commutative diagrams we use Xy-pic
\usepackage[all]{xy}

% We use 2cell for 2-commutative diagrams.
\xyoption{2cell}
% Theorem environments.
%
\theoremstyle{plain}
\newtheorem{theorem}[subsection]{Theorem}
\newtheorem{proposition}[subsection]{Proposition}
\newtheorem{lemma}[subsection]{Lemma}

\theoremstyle{definition}
\newtheorem{definition}[subsection]{Definition}
\newtheorem{example}[subsection]{Example}
\newtheorem{exercise}[subsection]{Exercise}
\newtheorem{situation}[subsection]{Situation}

\theoremstyle{remark}
\newtheorem{remark}[subsection]{Remark}
\newtheorem{remarks}[subsection]{Remarks}

\numberwithin{equation}{subsection}

% Macros
%
\def\lim{\mathop{\mathrm{lim}}\nolimits}
\def\colim{\mathop{\mathrm{colim}}\nolimits}
\def\Spec{\mathop{\mathrm{Spec}}}
\def\Hom{\mathop{\mathrm{Hom}}\nolimits}
\def\Ext{\mathop{\mathrm{Ext}}\nolimits}
\def\SheafHom{\mathop{\mathcal{H}\!\mathit{om}}\nolimits}
\def\SheafExt{\mathop{\mathcal{E}\!\mathit{xt}}\nolimits}
\def\Sch{\mathit{Sch}}
\def\Mor{\mathop{\mathrm{Mor}}\nolimits}
\def\Ob{\mathop{\mathrm{Ob}}\nolimits}
\def\Sh{\mathop{\mathit{Sh}}\nolimits}
\def\NL{\mathop{N\!L}\nolimits}
\def\CH{\mathop{\mathrm{CH}}\nolimits}
\def\proetale{{pro\text{-}\acute{e}tale}}
\def\etale{{\acute{e}tale}}
\def\QCoh{\mathit{QCoh}}
\def\Ker{\mathop{\mathrm{Ker}}}
\def\Im{\mathop{\mathrm{Im}}}
\def\Coker{\mathop{\mathrm{Coker}}}
\def\Coim{\mathop{\mathrm{Coim}}}

% Boxtimes
%
\DeclareMathSymbol{\boxtimes}{\mathbin}{AMSa}{"02}

%
% Macros for moduli stacks/spaces
%
\def\QCohstack{\mathcal{QC}\!\mathit{oh}}
\def\Cohstack{\mathcal{C}\!\mathit{oh}}
\def\Spacesstack{\mathcal{S}\!\mathit{paces}}
\def\Quotfunctor{\mathrm{Quot}}
\def\Hilbfunctor{\mathrm{Hilb}}
\def\Curvesstack{\mathcal{C}\!\mathit{urves}}
\def\Polarizedstack{\mathcal{P}\!\mathit{olarized}}
\def\Complexesstack{\mathcal{C}\!\mathit{omplexes}}
% \Pic is the operator that assigns to X its picard group, usage \Pic(X)
% \Picardstack_{X/B} denotes the Picard stack of X over B
% \Picardfunctor_{X/B} denotes the Picard functor of X over B
\def\Pic{\mathop{\mathrm{Pic}}\nolimits}
\def\Picardstack{\mathcal{P}\!\mathit{ic}}
\def\Picardfunctor{\mathrm{Pic}}
\def\Deformationcategory{\mathcal{D}\!\mathit{ef}}
\title{An ideal inside the Jacobson radical of smaller cohomological dimension pulls the cohomological dimension of a larger ideal down}
\author{The Stacks Project Authors}
\date{stacks-project:local-cohomology.tex:local-snapshot-20261012; GNU FDL 1.2 or later}
\begin{document}
\maketitle
\noindent\textit{Editorial note.} All delivered mathematical text is the author's own, copied from the pinned Stacks Project chapter source sources/stacks-local-cohomology/pinned-original-source.tex; no statement or proof step was rewritten, reordered or removed.\par\medskip
\noindent\textit{Editorial note.} Preamble provenance: the Stacks Project keeps one single shared preamble for every chapter; this card takes it directly from the staged sources/stacks-local-cohomology/preamble.tex as three verbatim spans (lines 8-12, 14-18 and 154-217), so no derived file was written and nothing was added or redefined. The wrapper class and the base packages are supplied by the assembler.\par\medskip
\noindent\textit{Editorial note.} Bibliography provenance: neither the counted statement nor its proof contains a citation (this lane rescans the whole span before writing the card), so the card carries no bibliography and no entry was reconstructed.\par\medskip
\noindent\textit{Editorial note.} Selection provenance: this card was opened under the file-based control contract (CONTROL.md), with the mechanical pre-card gate run on its statement and proof interval before the card was written.\par\medskip
\noindent\textit{Allowed standard prerequisites (editorial).}\begin{itemize}
\item The vanishing of local cohomology above the dimension of the support and the bound cd(A, I) \textbackslash{}leq dim(A). Author lemma lemma-cd-dimension, whose statement is reproduced verbatim as a context span of this card (source lines 758-764). Gives the a priori bound for the pair of ideals and the vanishing at the non-maximal prime.
\item The criterion for the comparison map between two local cohomology modules to be surjective in terms of the vanishing of the local cohomology at the primes of the difference. Author lemma lemma-isomorphism, whose statement is reproduced verbatim as a context span of this card (source lines 485-506). Reduces the vanishing needed for cd(A, I) < d to a surjectivity statement on the top local cohomology of the ring.
\end{itemize}\medskip
\begin{lemma}
\label{lemma-cd-top-vanishing}
Let $A$ be a Noetherian ring of dimension $d$. Let $I \subset I' \subset A$
be ideals. If $I'$ is contained in the Jacobson radical
of $A$ and $\text{cd}(A, I') < d$, then $\text{cd}(A, I) < d$.
\end{lemma}

\begin{lemma}
\label{lemma-cd-dimension}
Let $I \subset A$ be a finitely generated ideal of a ring $A$.
If $M$ is a finite $A$-module, then
$H^i_{V(I)}(M) = 0$ for $i > \dim(\text{Supp}(M))$.
In particular, we have $\text{cd}(A, I) \leq \dim(A)$.
\end{lemma}

\begin{lemma}
\label{lemma-isomorphism}
Let $I \subset I' \subset A$ be finitely generated ideals of a
Noetherian ring $A$. Let $M$ be an $A$-module. Let $i \geq 0$ be an integer.
Consider the map
$$
\Psi : H^i_{V(I')}(M) \to H^i_{V(I)}(M)
$$
The following are true:
\begin{enumerate}
\item if $H^i_{\mathfrak pA_\mathfrak p}(M_\mathfrak p) = 0$
for all $\mathfrak p \in V(I) \setminus V(I')$, then
$\Psi$ is surjective,
\item if $H^{i - 1}_{\mathfrak pA_\mathfrak p}(M_\mathfrak p) = 0$
for all $\mathfrak p \in V(I) \setminus V(I')$, then
$\Psi$ is injective,
\item if $H^i_{\mathfrak pA_\mathfrak p}(M_\mathfrak p) =
H^{i - 1}_{\mathfrak pA_\mathfrak p}(M_\mathfrak p) = 0$
for all $\mathfrak p \in V(I) \setminus V(I')$, then
$\Psi$ is an isomorphism.
\end{enumerate}
\end{lemma}

\begin{proof}
By Lemma \ref{lemma-cd-dimension} we know $\text{cd}(A, I) \leq d$.
We will use Lemma \ref{lemma-isomorphism} to show
$$
H^d_{V(I')}(A) \to H^d_{V(I)}(A)
$$
is surjective which will finish the proof. Pick
$\mathfrak p \in V(I) \setminus V(I')$. By our assumption
on $I'$ we see that $\mathfrak p$ is not a maximal ideal of $A$.
Hence $\dim(A_\mathfrak p) < d$. Then
$H^d_{\mathfrak pA_\mathfrak p}(A_\mathfrak p) = 0$
by Lemma \ref{lemma-cd-dimension}.
\end{proof}

\end{document}
